\documentclass[a4paper,12pt]{article}
\addtolength{\oddsidemargin}{-1.cm}
\addtolength{\textwidth}{2cm}
\addtolength{\topmargin}{-3cm}
\addtolength{\textheight}{3.5cm}
\makeindex


\usepackage[pdftex]{graphicx}
\usepackage{makeidx}
\usepackage{float}
\usepackage{hyperref}
\hypersetup{
    colorlinks=true,
    linkcolor=blue,
    filecolor=magenta,      
    urlcolor=cyan,
}


% define the title
\author{Team Chowers}
\title{Demo 1}
\begin{document}
\setlength{\parskip}{6pt}

% generates the title
	\begin{titlepage}
		\begin{center}
			\includegraphics[width=1\textwidth]{./up-logo.jpg}\\[1.5cm] 
			\textsc{\LARGE Department of Computer Science} \\ [.5cm]
			\textsc{\Large COS 301 - Atomic Game Engine} \\ [.5cm]
			\textsc{\Large Demo 1} \\ [.5cm]
			\line(1,0){450}\\[.5cm]
			\huge{\bfseries Software Requirements Specification and Project Planning}\\
			\line(1,0){450}\\[.5cm]
			\textsc{\LARGE Team Chowers}\\ [0.5cm]
			
			
			\textsc{\small Jocelyn Bondjobo (13232852)}\\
			\textsc{\small Sara Masilela 	(10126202)}\\
			\textsc{\small Keorapetse Shiko (12231992)}\\
			\textsc{\small Daniel Malangu (13315120)}\\
			
		\end{center}
	\end{titlepage}
	
\tableofcontents
\thispagestyle{empty}
\footnotesize
\normalsize

\newpage
\section{Introduction}
Atomic Game Engine is a relatively new platform for game development. It is open source and aims to compete with other engines like Unity and Unreal. We use Atomic for all of our game development at the moment and are looking to extend its capabilities.
The client, Johan van Staden from Luma Animation has requested the team to research the implementation to make the Atomic Game Engine work with the VR Technology as currently the Atomic Game Engine does not support any VR Technology. And, make sure ot runs at as high a framerate as possible. The team will also have to build an interactive experince using the new technology as a demo.

\newpage
\section{Vision}
Virtual reality is a going to be a big component of the hard-core game experience in the future and opening this technology up for open-source development is something that we would like to push for.
\newpage
\section{Background}

\newpage
\section{Software Architecture}
\subsection{Architecture Requirements}
\subsubsection{Architectural Scope}

\subsubsection{Access Channel Requirements}

\subsubsection{Quality Requirements}

 \begin{itemize}
 	\item[$\bullet$]Performance: 

 	\item[$\bullet$]Reliability:  

 	\item[$\bullet$]Scalability: 

 	\item[$\bullet$]Security:

 	\item[$\bullet$]Flexibility: 

	 \item[$\bullet$]Maintainability: 

 	\item[$\bullet$]Auditability:

 	\item[$\bullet$]Integrability:
 	\item[$\bullet$]Usability:

 	 \end{itemize}

\subsubsection{Integration Requirements}
\begin{flushleft}
\textbf{Integration Channels}

\textbf{Quality Requirements Integration}

\end{flushleft}

\subsubsection{Architectural Constraints}

\subsection{Architectural patterns or styles}
\subsection{Architectural tactics or strategies}

\newpage
\subsection{Use of reference architectures and frameworks}
\subsubsection{Examples of reference architectures}

\subsection{Access and Integration Channels}
\subsubsection{Access Channels}
\subsubsection{Integration Channels}

\newpage
\subsection{Technologies}	
	\subsubsection{Programming languages}
	
	\subsubsection{Operating Systems}
\newpage
\section{Project Planning and implementation}
	
\subsection{Required functionality}

\subsection{Process specifications}

\newpage
\section{Open Issues}

\subsection{References}

\end{document}
